\part{Unresolved comparisons and current gates}
\section{Statements not asserted as proved}

There is no unchecked theorem silently used in the preceding part.
The implications with hypotheses, including B8 and C15's universality
reduction, are proved implications; their hypotheses are not supplied
by this run. The following are the remaining identifications.

\paragraph{G1: the exact iterated character-independent functional.}
The common oscillator kernel, A1's normalized zero orbit, B11's
restriction seesaw, A5's finite-frame identity and A7's actual single
theta lift have been established. B17 now makes the actual wedge
a literal double lift $V\to P\to V$ by the genuine adjoint
kernel and a common positive inverse $Q^{-1}$.
It specializes the analytic last-step Bessel seesaw without a
new occurrence hypothesis. The operator is common to every
character for fixed spaces, kernel and levels.
The remaining \texttt{MainX} question is its stronger arithmetic
uniform relation-preservation statement. To apply it, one must
construct its actual common geometry,
tuples, archimedean maps, definite class kernels and arithmetic links,
and identify its normalized functional on the selected 032 vectors.
The concrete all-place tower fails the canonical node-gap and
signature requirements. The abstract two-step case is more flexible:
with global Witt index zero, one can take $m=2$ and the Bessel
unipotent rank zero. Its middle space must be totally definite.
A7's natural multiplicity space is totally definite precisely when
the selected first-list CM types agree; that direct realization then
gives a divisor-generated balanced line. B16 supplies a conditional
pole criterion for a different coherent definite source.
The arithmetic kernel/covolume and terminal-mass equations,
and an arithmetic interpretation of the analytic correction,
are the remaining matching problems. B17's upper space is
$A\perp\mathbb H$, of signature $(3,1)$ at every real place;
it is not a hidden definite inner form.

B19 narrows the arithmetic comparison to a finite CM-jet
functional on the actual polarized tori at a rational negative
line. Finite Heisenberg equivariance gives one theta-null
comparison scalar for every fixed-type finite input family;
the tangent periods, levels, stabilizers and Haar volumes are
explicit. The exact archive class-kernel identity can now be
tested on a finite basis after its polarized datum and frames
are matched. B20 supplies the literal isotropic doubled columns
for the exponent-one value using even integral reference sections.
It does not identify the resulting global automorphic functions
or provide high-order integral constant jets. Across different
balanced CM types, the theta-null determinant-period factors
still require a proof; balance of tangent lines alone is
insufficient to identify the full Hodge determinant.
A10 now constructs the common stabilized product theta module,
including the exact reference multipliers and a correspondence
proof of the balanced tangent-period relation. B21 identifies
the raw scalar obstruction as the specific ratio
$D_1D_2/(D_3D_4)$ of full Hodge period determinants. It proves
by an explicit complete-type calculation that this determinant
tensor need not balance. No transcendence conclusion follows;
the raw period relation, or its retained correction, and the
global weight/signature comparison remain explicit.

A11 settles the actual generic geometric host problem. Its
degree-at-most-four field extension is finite-unramified over
the real subfield, and its minimal CM-type padding, compatible
principalization and finite theta interpolation produce the
literal \texttt{Geo.Tor.Host/Datum/Obj}. The transported root
in \texttt{Obj.geom} accounts for every coefficient embedding.
The dimension cost is at most $64$ times the initial dimension.
This geometric host is different from the stronger global
\texttt{Formalism.BasicNotions.Host} used by \texttt{D7};
its additional Galois frame, arithmetic parameters and the
prescribed \texttt{StdPair} and tuple still need construction.
A16 proves that field-preserving enlargement cannot construct this
stronger host for arbitrary CM fields: its field is necessarily
abelian Galois, whereas the displayed nonnormal quartic CM field
and its nonabelian Galois normal closure cannot embed into it.
Thus the universal literal-host enlargement route is refuted;
the general comparison must use the field-independent proved
identities with their own arithmetic hypotheses.
The interpolation preserves algebraic first jets and records
the analytic comparison scalars, rather than identifying the
entire global theta family.

B23 gives an algebraic interpretation of A5's correction on
each full CM Hodge block. The non-Tate cup image of the actual
finite CM sources is an algebraic retract, constructed using
CM Hodge only on the source product. In a full multiplicity
and embedding block, the inverse is an algebraic correspondence
times $\nu_G|P_\chi|^2/|t_x|^2$. The common canonical-fibre
period $t_x$ is the specified Kodaira--Spencer determinant
divided by $(2\pi i)^2$. This removes the matrix-algebraicity
gate in that setting. The scalar and the raw theta-coordinate
comparisons remain explicit; the inverse of one fixed B17
integral kernel has not been identified with this operator.

A12 goes beyond abstract interpolation on the standard side.
For a specified compatible \texttt{StdPair}, finite
Heisenberg congruence projectors have actual
\texttt{IsThetaAt/IsRatDatum} witnesses. An explicit Gaussian
tail certificate makes their first-derivative matrix invertible;
the algebraic inverse retains the vacuum scalar, theta scalar,
common tangent period and isogeny matrices. When the prescribed
global \texttt{xiAt} has the stated support and period bounds,
its equality with the newly constructed array is exactly an
equality on one finite quotient. This does not supply the
stronger host, change the original tuple by definition, or
identify the global theta functions. The full jet value at
the origin is not prescribed by the derivative spanning result.

B24 identifies the actual B17 integral-kernel inverse in the
natural totally definite same-type sector. The source infinity
type is a character of the full compact group; its integral
reduces to finitely many rational masses. A common B19 CM
evaluation factor makes the kernel matrix projectively
algebraic and its image algebraically defined in the full CM
block. The actual inverse is consequently an explicit positive
scalar times an algebraic correspondence. A7 shows that this
natural sector is divisor-generated. For the general
exceptional balanced type the remaining operator condition is
projective algebraicity of the specified noncompact-source
Gram matrix, or a different coherent definite realization.

B25 removes the image-preservation problem in the finite-period
identity even for the actual noncompact middle source. Cayley
weak approximation and rational theta automorphy show that
finite adelic source translates span the same finite theta
image. Their common CM evaluation scalar makes that image
algebraic in coherent coordinates. The algebraic finite-frame
inverse can therefore preserve it. For every pair in this
image, including 032's actual fixed-input compact-place
variation, one has the exact identity
\[
 \langle f,g\rangle_G
   =s_\chi\sum_{j,\alpha}
       \overline{\mathcal P_{j,\alpha}(K_0f)}
                     \mathcal P_{j,\alpha}(g),
\]
with one algebraic correspondence $K_0$ and one named positive
scalar. Each corrected vector remains in the original B17
theta image, so the literal last-step seesaw applies there.
This gives the fully proved finite definite-period form of
the proposed comparison. The actual integral inverse $TT^*$
still has its own Gram-period matrix; the displayed finite
restriction inverse does not replace it, and the finite sum
is not asserted to be one prescribed arithmetic MainX tuple.

A13 now realizes any stipulated admissible finite tame character
family in a literal valid \texttt{TC} on an existing full
\texttt{Pre}. The coefficient enlargement preserves the entire
\texttt{Base}; the source constructor retains the chosen ambient
datum $J$, and the proved universal laws supply \texttt{Laws}
and \texttt{Law0}. With the unchanged \texttt{HasLetters},
\texttt{Later} and a same-$J$ arithmetic-jet field follow.
A14 supplies that specialized $p$-primary \texttt{TwDat} shape
for any prescribed finite local primary family at the stated
allowed split places. Starting from an existing full witness,
it reconstructs \texttt{ColW} while preserving the literal
\texttt{Pins}, \texttt{KsExact}, \texttt{FieldOK} and
\texttt{ValLL0}. The required coefficient-prime orbits are
created by a disjoint real multiquadratic extension split at $p$.
The actual arrays vanish at zero and have full prescribed rank
on their support. Their local incidence and determinant-character
spans, and the exact root-of-unity obstruction to global tame
characters, are proved. The remaining finite gate is the actual
Heisenberg/Fourier transform of these arrays followed by the
period-normalized first-derivative matrix. Raw array-window
selection alone is not the Gaussian-window construction of A12.
An arbitrary 032 normalization has not been asserted to satisfy
the literal \texttt{AmbIn.Ok} conditions.

B26 further reduces the actual noncompact-source integral Gram.
Its coefficient sections lie in the algebraic weight-three
source line, with the central CM determinant factor retained,
and their coherent coordinates have the common scalar
$\gamma=C_x/(t_x t_y)$. The Gram is an algebraic compression
of $\bigoplus_\lambda H_\lambda\otimes G_\lambda$, where
$G_\lambda$ is algebraic and $H_\lambda$ is the actual positive
multiplicity-period matrix. A single irreducible Hecke sector
therefore closes the actual $TT^*$ inverse up to one positive
scalar. Distinct sector norm ratios and higher multiplicity
matrices remain. Its cohomological source is
$H^s(Y,\bigotimes_\tau\mathbb V_\tau)$ with the recorded
constant line $\mathscr D=\mathscr C\otimes\bigotimes_\tau
\det\mathbb V_\tau$; neither the determinant comparison periods
nor the exact $4^{-s}$ measure factor is discarded.

B27 proves that the actual rank-two mixed source has
multiplicity-free finite Hecke sectors. The only nongeneric
rank-two parameter would be a determinant character at a
split finite place; strong approximation would then force
global derived-group invariance, contradicting the actual
weight-three discrete series. The proved generic pure-inner
multiplicity formula applies. Thus every $H_\lambda$ above
is a positive scalar $c_\lambda$. Each source-sector kernel
has a projectively algebraic inverse on its actual target
image, while the full kernel still has an algebraic
compression of these different scalar periods. B29 now proves
that their actual target images are orthogonal, by finite local
Howe uniqueness and the already fixed source infinity type.
The full inverse is therefore the exact sum
$(\nu_G|P_\chi|^2|\gamma|^2)^{-1}
\sum_\lambda c_\lambda^{-1}B_\lambda$ of algebraic sector
inverses with retained reciprocal scalar weights. Their exact
ratios remain. The operator also changes by exactly
$|c_t|^{-2}$ along the original compact-place variation.
B28's bound controls its norm without evaluating the source
periods.

C37 supplies the simultaneous analytic rank theorem for several
ordered rank-one differences at the same prime. It retains arbitrary
lattice representatives and all their moving character phases. A15 proves the exact ordered modulation identity and finite-block
injectivity from selfdual lattice eigencharacters and Fourier
inversion. A nonzero permitted array therefore supplies a nonzero
baseline in that block. The remaining application check is the
actual fixed-data and genuine-period-lattice transport as the
allowed archive prime grows; an arbitrary changing datum has
not been identified with a fixed positive section.

\paragraph{G2: turning the explicit certificates into a uniform
arithmetic lower bound for the prescribed iterated family.}
Q1--Q11 and B5--B7 give actual formulas under specified inputs.
The finite unipotent search yields $|J|\geq1/(4S_{\rm ext})$ with
an explicit grid size after its Gaussian envelope and lattice basis
are supplied. An arithmetic Gram bound needs a coefficient field,
denominator and bounds at all embeddings. Q9 provides resultant
denominators for specified integral Hecke matrices. Q11 transfers
the definite period-square certificate directly to the geometric
ball norm using explicit injective-chart constants; no unknown
ball Gram upper bound is needed. The remaining
step is the normalization and overlap data of the particular
iterated family in G1. Positivity alone cannot supply it, as B4
and B6 show. No universal cycle-degree estimate is asserted.
B19 now supplies the actual algebraic finite coefficient to
which the norm argument applies: its nonzero value has a lower
bound in terms of its denominator and its other conjugates,
multiplied by the specified common CM period and measure factor.
The arithmetic input bounds have not been supplied uniformly.
B22 removes the unknown absolute CM-period size by a specified
algebraic normalization: an even reference and an odd first jet
have explicit Gaussian-tail bounds, and the first-jet frame
makes the combined theta and tangent scalar exactly that
bounded analytic derivative. The resulting algebraic
denominators and conjugate bounds, and any loss of integral
saturation, must still be recorded. A9 also proves an explicit
constant-loss determinant derivative estimate under an integral
product-formal splitting, valid in every theta degree after
isotropic frame cancellation.

B28 supplies a further fixed-input certificate for the genuine
integral-kernel correction. On disjoint source polydisks,
weight-three integration gives
$\kappa_l=\nu_l(4\pi)^s\prod_j(\rho_{lj}^2-\rho_{lj}^4/2)$.
The actual source CM evaluation matrix is $(C_x/t_x)D$,
so $W\geq|C_x/t_x|^2(D^*\operatorname{diag}(\kappa_l)D)^{\mathsf t}$.
This bounds the least-norm source preimage and the actual
$TT^*$ inverse. The extra source-fibre period cancels exactly.
The remaining numerical data are the specified finite
algebraic minor and its denominators/conjugates, chart radii,
original measures and the retained target theta factor.
No uniform estimate for unspecified arithmetic data is supplied.

\paragraph{G3: an exceptional non-CM Hodge summand.}
C16 supplies infinitely many non-CM elliptic source factors, and
C20 adds a geometrically simple genus-two Jacobian with
endomorphism ring $\Z$, on a ball cover of degree at most1764.
C24 gives non-CM balanced $\Q(i)$-Weil sources in every
even dimension. C25 specifies a geometrically simple fourfold
with endomorphism ring $\Z[i]$, N\'eron--Severi rank one,
and entirely exceptional Weil space; its ball cover has degree
at most43264. Thus the source and the exceptional target
are both concrete.
C13 gives an explicit algebraic projector onto a Weil space.
C32 now supplies the nonzero projected cycle for the named
C25 target. Its integral class is pure Weil and has exact
self-intersection $73728$. Pullback by $[2+i]$ supplies the
second rational direction, so the exceptional Weil-space
extraction is complete for this target and the displayed
cyclic-four family.
The known curve-product maps in C12, C16 and C20 are subject to the
same C14 vanishing whenever the target has the generic Weil-type
divisor algebra. C32 uses a signed split-plane cycle, refined
pullback and relative hyperplane sections on the fourth power
of the actual ball cover; its image has self-intersection
$d^8\cdot73728$. The remaining task is to extend this
higher-correspondence construction to other non-CM Hodge
summands, or to identify a literal single-ball theta realization.
C15's characteristic-zero universality implication is a source
reduction, not a proof of that conjecture or of the extraction gate.
C27 makes the map-based version a finite quartic cup-product
test on the $A$-isotypic Albanese homomorphism space. Pulling
back existing zero pairings to another finite cover cannot
repair them; genuinely new maps and mixed cup coefficients
are required. A fixed generating source map also has an
injective target Kodaira--Spencer obstruction, so a deformation
argument must use a higher relative cycle instead of retaining
that degree-one quotient.
The countability result restricts fixed arithmetic towers over
$\C$; it does not preclude the dense algebraic families constructed
here or a further deformation argument.

C28 now excludes the entire single-fibration base-change
tower for the specified C25 target, including all previously
unknown maps on those covers. The actual geometric generic
fibre has genus $76$; its Jacobian splits into factors of
dimensions two and four carrying $\Q(\zeta_5)$. It has no
homomorphism to the simple fourfold with endomorphism algebra
$\Q(i)$. Consequently every map on every such base change
factors through its Stein base curve and the full C27 quartic
vanishes. A map-based solution must therefore change the
generic fibre, or be replaced by a higher correspondence.

C30 also excludes the natural simultaneous two-Fermat-base
construction. Its finite cyclic model has no additional
coherent first cohomology, and resolving its explicit Fermat
cones adds only CM Jacobian quotients. The Albanese is
isogenous to the product of the two base Jacobians and these
CM factors. Thus every map to C25 factors, up to multiplication,
through the product of the two curves; C14 again makes the
entire quartic zero. An arbitrary finite cover need not arise
this way, and remains outside this calculation.

C32 also constructs integral pure Weil cycles on every C24
cyclic-four Jacobian of dimension $2n$, with exact square
$2(-1)^n4^{2n-1}((2n)!)^2$. Its $[2+i]$ pullback spans
the other rational direction, and the actual ball power has
image multiplier $d^{2n}$. The remaining scope is other
non-CM Hodge summands and the additional N\'eron--Severi
condition needed to call every higher-dimensional member
exceptional outside the divisor algebra.

The C28 and C30 exclusions therefore remain precise statements
about maps in those towers, not barriers to the constructed
higher correspondence. C32 uses Schoen's paired-branch geometry
and checks the order-four exceptional inertia directly. It
does not import Markman's cycle. The supplied fourfold theorem
already asserts algebraicity; the new result here is the
specified cycle, its normalization and its actual ball-power
realization.

\paragraph{G4: the chosen integral correspondence and connection.}
C17 identifies the exact horizontal correspondence and local-model
data needed to compare the 032 image cycle with a Kudla--Rapoport
special cycle. C06--C08 prevent replacing those data by an
undifferentiated positivity claim. C11 proves the sharp integral
height-two derivative estimate. C19 removes its artificial
polynomial-coordinate hypothesis and proves integrality of the
CM de Rham normalization coefficient $e_2^*$.
A8 identifies the actual completed 032 tensor with the matched
canonical harmonic jet, including its embedding labels and
factorials. On a matched CM elliptic factor, A8/C19 give the
actual reduced-theta derivative bounds and their prime-to-$p$
level normalization. C21 proves the actual product decomposition
at primes split in $F$ and inert in $E$, and the strong
mixed-derivative bound for canonically matched equal-degree
ratios with integral sections and unit denominator.
For arbitrary integer tensor degree the elliptic bound has
the weaker positive slope, which C22 proves optimal through
exact infinite subsequences on two actual tame theta sections.
Thus transferring the ratio slope to all individual sections
is decided negatively in this normalization.
A9 resolves the positive-degree frame problem for the precise
archive isotropic columns: all higher canonical gauge terms
cancel to an integral unit, and the ordinary formal depth
equals the actual local theta depth. C26 proves an independent
actual elliptic two-column instance. B20 completes the first
determinant value to these columns under explicit integral
evaluation and saturation hypotheses. These conclusions do
not extend the raw one-copy bound refuted in C22.

C23 constructs the actual A6 cycle over a specified good model
and specializes its nonzero class. Its pure CM kernels are
finite combinations of normalized Poincar\'e biextensions,
the exact category of linear special homomorphisms in SGA~7~I,
VIII~3.1--3.2. C29 and C31 realize these kernels by signed
labelled homomorphism-lifting graphs on actual PEL boundary
charts, for elliptic and higher CM factors respectively.
C29 retains the exact integral multiplier $24$ or $48$;
C31's higher-CM construction instead uses the specified
prime-to-$p$ isogeny chart and zero kernel labels.

C33 proves that these graphs are exactly the labelled horizontal
interior special-cycle closures on the stated formal charts.
The good source has zero period lattice, so Lan's fully faithful
theorem for arbitrary homomorphisms identifies the extension
condition with $h^\vee c=0$. The three positive Hermitian norms,
characteristic-zero deformation functor and saturated scheme
structure are explicit. Refined boundary restriction gives the
graph before the proper polarization contraction on its
abelian parameter space.

C34 supplies the relative correspondence to the original compact
ball on the literal product of its good model and the actual
small PEL hosts. Both reference extension and dual-period
coordinates are zero, and the degeneration splits as stated.
Pulling back the exact normalized line bundles along the actual
ball maps and applying the proper Chern-class contraction gives
$\mathcal C'=\delta\mathcal C$ for the chosen original cycle.
The finite cover degree, including any $p$-divisibility, and all
CM projector coefficients remain explicit. Thus the chosen
horizontal-cycle comparison is proved at these good models.
The remaining quantitative comparison concerns the exact theta
normalizations and integral denominators. Vertical local
Rapoport--Zink cycles and their derived intersection numbers
require their own equations; they are not inferred from this
horizontal identity. No proper projection of the full toroidal
formal chart is used.

C35 fixes an exact local quantitative datum for those same
horizontal closures. Their Hermitian Gram is integrally
congruent to $\operatorname{diag}(a\bar a,1,1)$; their derived
intersection has no higher Tor, and its refined boundary
restriction is $B[a]$ with degree $N_{E/\Q}(a)$.
At the specified inert factors its identity length is
$p^{2\sum r_\nu}$. Its coherent density polynomial is also
computed with the same Gram. Its formal derivative is not
identified with a vertical arithmetic intersection.

These gates are recorded with dated changes in \texttt{LEDGER.md}.
The mandated three-hour progress review has not been replaced by
this mathematical inventory. Its timing is controlled by the actual
run clock and recorded separately.
